%% TRANSCRIBED (PDF 55-77 = printed II-6..II-28)
%% Sec. 2 "Construction of T_m and S_m".  The top of PDF 55 (II-6) up to
%% ``What happens when Spec(k') is not connected?''  belongs to the last
%% Exercise of Sec. 1.3 and is transcribed in ch2_s1.tex.
%% NOTE: the final paragraph of 2.8 spills onto II-29 (PDF 78), just before
%% the Sec. 3 heading; it is transcribed at the end of this file, so that
%% ch2_s3.tex can start directly at \sechead{3}.
%% NOTE (subscript "o"): the subscripts of $V_o$, $\phi_o$, $\End_K(A)_o$,
%% $\End(\tilde{A}_v)_o$ and of Weil's type $(A_o)$ are printed as a small
%% round letter ``o'' (same convention as $K_o$ in chap. I, 1.2 and $A_o$
%% in \S\ref{sec:II.1.3}); transcribed literally as ``o''.  (Chap. III,
%% 2.2-2.3 denotes the same objects $V_0$, $\phi_0$.)

\sechead{2}{Construction of $T_m$ and $S_m$}{CONSTRUCTION OF $T_m$ AND $S_m$}

\subsechead{2.1}{Idèles and idele-classes}

We defined in chap.~\ref{chap:I}, \S\ref{sec:I.2.1} the set $\Sigma_K$ of
finite places of the number field $K$.  Let now $\Sigma_K^\infty$ be the
set of equivalence classes of archimedian absolute values of $K$, and let
$\bar{\Sigma}_K$ be the union of $\Sigma_K$ and $\Sigma_K^\infty$.  If
$v \in \bar{\Sigma}_K$ then $K_v$ denotes the \emph{completion} of $K$
with respect to $v$.  For $v \in \Sigma_K^\infty$ we have $K_v = \R$ or
$K_v = \C$, and $K_v$ is ultrametric if $v \in \Sigma_K$.  For
$v \in \Sigma_K$, the group of units of $K_v$ is denoted by $U_v$.  The
\emph{idèle group} $I$ of $K$ is the subgroup of
$\prod\limits_{v \in \bar{\Sigma}_K} K_v^*$ consisting of the families
$(a_v)$ with $a_v \in U_v$, for almost all $v$; it is given a topology by
decreeing that the subgroup (with the product topology)
\[
\prod_{v \in \Sigma_K^\infty} K_v^* \times \prod_{v \in \Sigma_K} U_v
\]
be open.  We embed $K^*$ into $I$ by sending $a \in K^*$ onto the idèle
$(a_v)$, where $a_v = a$ for all $v$.  The topology induced on $K^*$ is
the discrete topology.  The quotient group $C = I/K^*$ is called the
\emph{idèle-class group} of $K$.  (For all this, see
\cite{cassels-frohlich}, \cite{lang-abelian}, or \cite{weil-bnt}.)

Let $S$ be a finite subset of $\Sigma_K$.  Then by a \emph{modulus of
support} $S$ we mean a family $m = (m_v)_{v \in S}$ where the $m_v$ are
integers $\geq 1$.

If $v \in \bar{\Sigma}_K$ and $m$ is a modulus of support $S$, we let
$U_{v,m}$ denote the connected component of $K_v^*$ if
$v \in \Sigma_K^\infty$, the subgroup of $U_v$ consisting of those
$u \in U_v$ for which $v(1-u) \geq m_v$ if $v \in S$, and $U_v$ if
$v \in \Sigma_K - S$.  The group $U_m = \prod\limits_v U_{v,m}$ is
\emph{an open subgroup} of $I$.

If $E$ is the group of units of $K$, let $E_m = E \cap U_m$.  The
subgroup $E_m$ is of finite index in $E$.  (Conversely, by a theorem of
Chevalley (\cite{chevalley}, see also \cite{serre-congr}, n°~3.5) every
subgroup of finite index in $E$ contains an $E_m$ for a suitable modulus
$m$.)

Let $I_m$ be the quotient $I/U_m$ and $C_m$ the quotient
$I/K^*U_m = C/(\text{Image of } U_m \text{ in } C)$.  One then has the
exact sequence
\[
1 \longrightarrow K^*/E_m \longrightarrow I_m \longrightarrow C_m
\longrightarrow 1.
\]

\emph{The group} $C_m$ \emph{is finite}; in fact, the image of $U_m$ in
$C$ is open, hence contains the connected component $D$ of $C$, and the
group $C/D$ is known to be compact (see \cite{lang-abelian},
\cite{weil-bnt}).  Moreover, any open subgroup of $I$ contains one of the
$U_m$'s, hence $C/D$ \emph{is the projective limit} of the $C_m$'s.
Class field theory (cf.  for instance \cite{cassels-frohlich}), gives
\emph{an isomorphism of} $C/D = \varprojlim\, C_m$ \emph{onto the Galois
group} $G^{\mathrm{ab}}$ of the maximal abelian extension of $K$.

\runhead{Remark}

A more classical definition of $C_m$ is as follows.  Let
$\mathrm{Id}_S$ be the group of fractional ideals of $K$ prime to $S$,
and $P_{S,m}$ the subgroup of principal ideals $(\gamma)$, where
$\gamma$ is totally positive and $\gamma \equiv 1 \bmod.\, m$ (i.e.
$\gamma$ belongs to $U_{v,m}$ for all $v \in S$ and
$v \in \Sigma_K^\infty$).  Let $\mathrm{Cl}_m = \mathrm{Id}_S/P_{S,m}$.
We have the exact sequence:
\[
1 \longrightarrow P_{S,m} \longrightarrow \mathrm{Id}_S
\longrightarrow \mathrm{Cl}_m \longrightarrow 1.
\]

For each $\mathbf{a} = \prod\limits_{v \notin S} v^{a_v}
\in \mathrm{Id}_S$, choose an idèle $\alpha = (\alpha_v)$, with
$\alpha_v \in U_{v,m}$ if $v \in S$ or $v \in \Sigma_K^\infty$, and
$v(\alpha_v) = a_v$ if $v \in \Sigma_K - S$.  The image of $\alpha$ in
$I_m = I/U_m$ depends only on $\mathbf{a}$.  We then get a homomorphism
$g\colon \mathrm{Id}_S \to I_m$.  One checks readily that $g$ extends to
a commutative diagram
\[
\begin{array}{ccccccc}
1 \to & P_{S,m} & \to & \mathrm{Id}_S & \to & \mathrm{Cl}_m & \to 1\\
      & \downarrow & & \downarrow\, g & & \downarrow\, f & \\
1 \to & K^*/E_m & \to & I_m & \to & C_m & \to 1\,,
\end{array}
\]
and that $f\colon \mathrm{Cl}_m \to C_m$ \emph{is an isomorphism}; hence
$C_m$ can be identified with the ideal class group mod $m$ (and this
shows again that it is finite).

\subsechead{2.2}{The groups $T_m$ and $S_m$}

We are now in a position to apply the group construction of
\S\ref{sec:II.1.3}.  We take for exact sequence
\eqref{eq:II.1.3.star} the sequence
\[
1 \longrightarrow K^*/E_m \longrightarrow I_m \longrightarrow C_m
\longrightarrow 1
\]
and for $A$ the algebraic group $T_m = T/\bar{E}_m$, where $E_m$ is as
before, $T$ is the torus $\mathrm{R}_{K/\Q}(\mathbf{G}_{m/K})$ defined in
\S\ref{sec:II.1.1}, and $\bar{E}_m$ is the Zariski closure of $E_m$ in
$T$, cf.~\S\ref{sec:II.1.2}.

The construction of \S\ref{sec:II.1.3} now yields a
\emph{$\Q$-algebraic group} $S_m$ with an algebraic morphism
$T_m \to S_m$ and a group homomorphism $\epsilon\colon I_m \to S_m(\Q)$.
The sequence
\[
1 \longrightarrow T_m \longrightarrow S_m \longrightarrow C_m
\longrightarrow 1
\]
is exact ($C_m$ being identified with the corresponding constant
algebraic group) and the diagram
\begin{equation}
\tag{**}\label{eq:II.2.2.starstar}
\begin{array}{ccccccc}
1 \to & K^*/E_m & \to & I_m & \to & C_m & \to 1\\
      & \downarrow & & \downarrow\, \epsilon & & \downarrow\, \mathrm{id.} & \\
1 \to & T_m(\Q) & \to & S_m(\Q) & \to & C_m & \to 1
\end{array}
\end{equation}
is commutative.

\runhead{Remark}

Let $m'$ be another modulus; assume $m' \geq m$, i.e.
$\Supp(m') \supset \Supp(m)$ and $m'_v \geq m_v$ if $v \in \Supp(m)$.
From the inclusion $U_{m'} \subset U_m$ one deduces maps
$T_{m'} \to T_m$ and $I_{m'} \to I_m$, whence a morphism
$S_{m'} \to S_m$.  Hence the $S_m$'s form a projective system; their
limit is a proalgebraic group over $\Q$, extension of the profinite group
$C/D = \varprojlim\, C_m$ by a torus.

\begin{exercises}
\begin{nlist}
\item Let $\bar{E}_m(\Q)$ be the Zariski-closure of $E_m$ in
$K^* = T(\Q)$.  Show that the kernel of
$\epsilon_m\colon I/U_m \to S_m(\Q)$ is the image of
$\bar{E}_m(\Q) \to I/U_m$.

\item Let $H_{m'/m}$ be the kernel of $S_{m'} \to S_m$, where
$m' \geq m$.
\begin{alist}[label=\alph*)]
\item Show that $H_{m'/m}$ is a finite subgroup of $(S_{m'}(\Q))$ and
that it is contained in the image of $\epsilon_{m'}$.
\item Construct an exact sequence (cf.~Exer.~1)
\[
1 \longrightarrow (E_m \cap \bar{E}_{m'}(\Q))/E_{m'}
\longrightarrow U_m/U_{m'} \longrightarrow H_{m'/m} \longrightarrow 1.
\]
\end{alist}
\end{nlist}
\end{exercises}

\subsechead{2.3}{The canonical \textit{l}-adic representation with values in $S_m$}

Let $m$ be a modulus, and let $\ell$ be a prime number.  Let
$\epsilon\colon I \to I_m \to S_m(\Q)$ be the homomorphism defined in
\S\ref{sec:II.2.2}.  Let $\pi\colon T \to S_m$ be the algebraic morphism
$T \to T_m \to S_m$; by taking points with values in $\Q_\ell$, $\pi$
defines a homomorphism
\[
\pi_\ell\colon T(\Q_\ell) \longrightarrow S_m(\Q_\ell).
\]

Since $K \tens \Q_\ell = \prod\limits_{v|\ell} K_v$, the group
$T(\Q_\ell)$ can be identified with
$K_\ell^* = \prod\limits_{v|\ell} K_v^*$, and is therefore a direct
factor of the idèle group $I$.  Let $\mathrm{pr}_\ell$ denote the
projection of $I$ onto this factor.  The map
\[
\alpha_\ell = \pi_\ell \circ \mathrm{pr}_\ell\colon I
\longrightarrow T(\Q_\ell) \longrightarrow S_m(\Q_\ell)
\]
is a continuous homomorphism.

\begin{lemma}
$\alpha_\ell$ and $\epsilon$ coincide on $K^*$.
\end{lemma}

This is trivial from the commutativity of the diagram
\eqref{eq:II.2.2.starstar} of \S\ref{sec:II.2.2}.

Now, let $\epsilon_\ell\colon I \to S_m(\Q_\ell)$ be defined by
\begin{equation}
\tag{***}\label{eq:II.2.3.starstarstar}
\begin{gathered}
\epsilon_\ell(a) = \epsilon(a)\,\alpha_\ell(a^{-1})\\
\text{i.e.}\quad \epsilon_\ell = \epsilon \cdot \alpha_\ell^{-1}
\end{gathered}
\end{equation}
(If $a \in I$, write $a_\ell$ the $\ell$-component of $a$.  Then
\[
\epsilon_\ell(a) = \epsilon(a)\,\pi_\ell(a_\ell^{-1}).)
\]

By the lemma, $\epsilon_\ell$ is trivial on $K^*$ and, hence, defines a
map $C \to S_m(\Q_\ell)$; since $S_m(\Q_\ell)$ is totally disconnected
(it is an $\ell$-adic Lie group), the latter homomorphism is trivial on
the connected component $D$ of $C$.  We have already recalled that $C/D$
may be identified with the Galois group $G^{\mathrm{ab}}$ of the maximal
abelian extension of $K$.  So we end up with a homomorphism
$\epsilon_\ell\colon G^{\mathrm{ab}} \to S_m(\Q_\ell)$, i.e.  with an
\emph{\textit{l}-adic representation of} $K$ \emph{with values in} $S_m$
(cf.~chap.~\ref{chap:I}, \S\ref{sec:I.2.3}).

This representation is \emph{rational} in the sense of
chap.~\ref{chap:I}, \S\ref{sec:I.2.3}.  More precisely, let
$v \notin \Supp(m)$, and let $f_v \in I$ be an idèle which is a
uniformizing parameter at $v$, and which is equal to $1$ everywhere
else; let $F_v = \epsilon(f_v)$ be the image of $f_v$ in $S_m(\Q)$.
With these notations we have:

\begin{proposition}
\begin{alist}[label=\alph*)]
\item The representation $\epsilon_\ell\colon G^{\mathrm{ab}} \to
S_m(\Q_\ell)$ is a rational representation with values in $S_m$.
\item $\epsilon_\ell$ is unramified outside $\Supp(m) \cup S_\ell$,
where $S_\ell = \{v\,|\,p_v = \ell\}$.
\item If $v \notin \Supp(m) \cup S_\ell$, then the Frobenius element
$F_{v,\epsilon_\ell}$ (cf.~chap.~\ref{chap:I}, \S\ref{sec:I.2.3}) is
equal to $F_v \in S_m(\Q)$.
\end{alist}
\end{proposition}

\begin{proof}
It is known that the class field isomorphism
$C/D \xrightarrow{\ \sim\ } G^{\mathrm{ab}}$ maps $K_v^*$ (resp.  $U_v$)
onto a dense subgroup of the decomposition group of $v$ in
$G^{\mathrm{ab}}$ (resp.  onto the inertia group of $v$ in
$G^{\mathrm{ab}}$), and that a uniformizing element $f_v$ of $K_v^*$ is
mapped onto the Frobenius class of $v$.

If $v \notin \Supp(m)$ and $a \in U_v$, then $\epsilon(a) = 1$; if
moreover $p_v \neq \ell$, $\alpha_\ell(a) = 1$, hence
$\epsilon_\ell(a) = 1$ and $\epsilon_\ell$ is unramified at $v$; this
proves \emph{b}).  For such a $v$, we have
$\epsilon_\ell(f_v) = \epsilon(f_v) = F_v$; hence \emph{c}), and
\emph{a}) follows from \emph{c}).
\end{proof}

\begin{corollary}
The representations $\epsilon_\ell$ form a system of strictly compatible
\textit{l}-adic representations with values in $S_m$.
\end{corollary}

We also see that the exceptional set of this system is contained in
$\Supp(m)$; for an example where it is different from $\Supp(m)$, see
Exercise~2.

\runhead{Remark}

By construction, $\epsilon_\ell\colon I \to S_m(\Q_\ell)$ is given by
$x \to \pi_\ell(x^{-1})$ on the open subgroup
$U_{\ell,m} = \prod\limits_{v|\ell} U_{v,m}$ of $K_\ell^*$.  Hence,
$\Ima(\epsilon_\ell)$ contains
$\pi_\ell(U_{\ell,m}) \subset T_m(\Q_\ell) \subset S_m(\Q_\ell)$, and is
an \emph{open subgroup} of $S_m(\Q_\ell)$.  This open subgroup maps onto
$C_m$, as remarked above.  These properties imply, in particular, that
$\Ima(\epsilon_\ell)$ is \emph{Zariski-dense} in $S_m$.

\begin{exercises}
\begin{nlist}
\item[(1)] Let $K = \Q$, $\Supp(m) = \emptyset$.
\begin{alist}[label=\alph*)]
\item Show that $E_m = \{1\}$, $C_m = \{1\}$, hence $T_m = S_m = \Gm$
and $S_m(\Q) = \Q^*$, $S_m(\Q_\ell) = \Q_\ell^*$.
\item Show that $I$ is the direct product of its subgroups $I_m$ and
$\Q^*$; hence any $a \in I$ may be written as
\[
a = u \cdot \gamma \qquad u \in U_m,\ \gamma \in \Q^*\,.
\]
Show that, if $a = (a_p)$, one has
\[
\epsilon(a) = \gamma = \mathrm{sgn}(a_\infty)
\prod_p\, p^{v_p(a_p)}\,.
\]
\item Show that
\[
\rho_\ell(a) = \gamma \cdot a_\ell^{-1},
\]
and
\[
F_p = p.
\]
\item Show that $\rho_\ell$ coincides with the character $\chi_\ell$ of
chap.~\ref{chap:I}, \S\ref{sec:I.1.2}.
\end{alist}

\item[(2)] Let $K = \Q$, $\Supp(m) = \{2\}$ and $m_2 = 1$.  Show that
the groups $E_m$, $C_m$, $T_m$, $S_m$ coincide with those of
Exercise~1, hence that the exceptional set of the corresponding system
is empty.
\end{nlist}
\end{exercises}

\subsechead{2.4}{Linear representations of $S_m$}

We recall first some well known facts on representations.

a) Let $k$ be a field of characteristic $0$; let $H$ be an affine
commutative algebraic group over $k$.  Let
$X(H) = \Hom_{\bar{k}}(H_{/\bar{k}}, \mathbf{G}_{m/\bar{k}})$ be the
group of characters of $H$ (of degree $1$).  Here we write the
characters of $X(H)$ multiplicatively.  The group
$G = \Gal(\bar{k}/k)$ acts on $X(H)$.

Let $\Lambda$ be the affine algebra of $H$, and let
$\bar{\Lambda} = \Lambda \tens_k \bar{k}$ be the one of $H_{/\bar{k}}$.
Every element $\chi \in X(H)$ can be identified with an invertible
element of $\bar{\Lambda}$.  Hence, by linearity, a homomorphism
\[
\alpha\colon \bar{k}[X(H)] \longrightarrow \bar{\Lambda}
\]
where $\bar{k}[X(H)]$ is the group algebra of $X(H)$ over $\bar{k}$.
This is a $G$-homomorphism if the action of $G$ is defined by
$s(\sum_\chi a_\chi\,\chi) = \sum_\chi s(a_\chi)s(\chi)$ for
$a_\chi \in \bar{k}$ and $\chi \in X(H)$.  It is well-known (linear
independence of characters) that $\alpha$ is \emph{injective}.  It is
\emph{bijective} if and only if $H$ is a group of multiplicative type
(cf.~\S\ref{sec:II.1.3}, remark~2).  Hence we may identify
$\bar{k}[X(H)]$ with a subalgebra of $\bar{\Lambda}$.

b) Let $V$ be a finite-dimensional $k$-vector space and let
\[
\phi\colon H \longrightarrow \GLV
\]
be a \emph{linear representation} of $H$ into $V$.  Assume $\phi$ is
\emph{semi-simple} (this is always the case if $H$ is of multiplicative
type).  We associate to $\phi$ its \emph{trace}
\[
\theta_\phi = \sum_\chi n_\chi(\phi)\,\chi
\]
in $\Z[X(H)]$, where $n_\chi(\phi)$ is the multiplicity of $\chi$ in the
decomposition of $\chi$ over $\bar{k}$.
%% NOTE (II-14): "the multiplicity of $\chi$ in the decomposition of
%% $\chi$ over $\bar k$" is as printed (one would expect "$\phi$" in
%% second place); kept as printed.

We have $\theta_\phi(h) = \mathrm{Tr}(\phi(h))$ for any point $h$ of $H$
(with value in any commutative $k$-algebra).  Let
$\mathrm{Rep}_k(H)$ be the set of isomorphism classes of linear
semi-simple representations of $H$.  If $k_1$ is an extension of $k$,
then scalar extension from $k$ to $k_1$ defines a map
$\mathrm{Rep}_k(H) \to \mathrm{Rep}_{k_1}(H_{/k_1})$ which is easily
seen to be \emph{injective}.  We say that an element of
$\mathrm{Rep}_{k_1}(H_{/k_1})$ \emph{can be defined over} $k$, if it is
in the image of this map.

\begin{proposition}[1]\label{prop:II.2.4.1}
The map $\phi \mapsto \theta_\phi$ defines a bijection between
$\mathrm{Rep}_k(H)$ and the set of elements
$\theta = \sum_\chi n_\chi\,\chi$ of $\Z[X(H)]$ which satisfy:
\begin{alist}
\item $\theta$ is invariant by $G$ (i.e.  $n_\chi = n_{s(\chi)}$ for all
$s \in G$, $\chi \in X(H))$.
\item $n_\chi \geq 0$ for every $\chi \in X(H)$.
\end{alist}
\end{proposition}

\begin{proof}
The injectivity of the map $\phi \mapsto \theta_\phi$ is well-known (and
does not depend on the commutativity of $H$).  To prove surjectivity,
consider first the case where $\theta$ has the form
$\theta = \sum \chi^{(i)}$ where $\chi^{(i)}$ is a full set of different
conjugates of a character $\chi \in X(H)$.  If $G(\chi)$ is the subgroup
of $G$ fixing $\chi$, then
\begin{equation}
\tag{$*$}\label{eq:II.2.4.star}
\theta = \sum_{s \in G/G(\chi)} s(\chi)\,.
\end{equation}

The fixed field $k_\chi$ of $G(\chi)$ in $\bar{k}$ is the smallest
subfield of $\bar{k}$ such that $\chi \in \Lambda \tens_k k_\chi$.
Consider $\chi$ as a representation of degree $1$ of $H_{/k_\chi}$.
One gets, by restriction of scalars to $k$, a representation $\phi$ of
$H$ of degree $[k_\chi : k]$.  One sees easily that the trace
$\theta_\phi$ of $\phi$ is equal to $\theta$.  The surjectivity of
$\phi \mapsto \theta_\phi$ now follows from the fact that any $\theta$
satisfying (a) and (b) is a sum of elements of the form
\eqref{eq:II.2.4.star} above.
\end{proof}

\begin{corollary}
In order that $\phi_1 \in \mathrm{Rep}_{k_1}(H_{/k_1})$ can be defined
over $k$, it is necessary and sufficient that
$\theta_{\phi_1} \in \Lambda \tens_k k_1$ \emph{belongs to} $\Lambda$.
\end{corollary}

(c) We return now to the groups $S_m$:

\begin{proposition}[2]\label{prop:II.2.4.2}
Let $k_1$ be an extension of $k$ and let
$\phi \in \mathrm{Rep}_{k_1}(S_{m/k_1})$.  The following properties are
equivalent:
\begin{rlist}
\item $\phi$ can be defined over $k$,
\item For every $v \notin \Supp(m)$, the coefficients of the
characteristic polynomial $\phi(F_v)$ belong to $k$,
\item There exists a set $\Sigma$ of places of $k$ of density $1$
(cf.~chap.~\ref{chap:I}, \S\ref{sec:I.2.2}) such that
$\mathrm{Tr}(\phi(F_v)) \in k$ for all $v \in \Sigma$.
\end{rlist}
\end{proposition}

\begin{proof}
The implications (i) $\Rightarrow$ (ii) $\Rightarrow$ (iii) are trivial.
To prove (iii) $\Rightarrow$ (i) we need the following lemma.
\end{proof}

\begin{lemma}
The set of Frobeniuses $F_v$, $v \in \Sigma$, is dense in $S_m$ for the
Zariski topology.
\end{lemma}

\begin{proof}
Let $X$ be the set of all $F_v$'s, $v \in \Sigma$, and let $\ell$ be a
prime number.  Let $\bar{X} \subset S_m$ (resp.  $\bar{X}_\ell \subset
S_m(\Q_\ell)$) the closure of $X$ in the Zariski topology (resp.
$\ell$-adic topology).  It is clear that
$\bar{X}_\ell \subset \bar{X}(\Q_\ell)$.  On the other hand,
\v{C}ebotarev's theorem (cf.~chap.~\ref{chap:I}, \S\ref{sec:I.2.2})
implies that $\bar{X}_\ell = \Ima(\epsilon_\ell)$
(cf.~\S\ref{sec:II.2.3}).  The set $\Ima(\epsilon_\ell)$, however, is
Zariski dense in $S_m$ (cf.~Remark in \S\ref{sec:II.2.3}).  Hence
$\bar{X} = S_m$, which proves the lemma.
\end{proof}

Let us now prove that (iii) $\Rightarrow$ (i).  Let $\theta_\phi$ be the
trace of $\phi$ in $\Lambda \tens_k k_1$, where $\Lambda$ is the affine
algebra of $H = S_{m/k}$.  Let $\{\ell_\alpha\}$ be a basis of the
$k$-vector space $k_1$, with $\ell_{\alpha_0} = 1$ for some index
$\alpha_0$.  We have
$\theta_\phi = \sum_\alpha \lambda_\alpha \tens \ell_\alpha$
($\lambda_\alpha \in \Lambda$); hence
$\mathrm{Tr}(\phi(h)) = \theta_\phi(h) = \sum_\alpha
\lambda_\alpha(h)\,\ell_\alpha$ for all $h \in H(k_1)$.  Take
$h = F_v$, with $v \in \Sigma$.  Since $F_v$ belongs to $H(k)$ we have
$\lambda_\alpha(F_v) \in k$ for all $\alpha$; since
$\mathrm{Tr}(\phi(F_v)) \in k$, we get $\lambda_\alpha(F_v) = 0$ for all
$\alpha \neq \alpha_0$.  By the lemma, the $F_v$'s, $v \in \Sigma$, are
Zariski-dense in $H$; hence $\lambda_\alpha = 0$ for $\alpha \neq
\alpha_0$ and $\theta_\phi = \lambda_{\alpha_0}$ belongs to $\Lambda$
and (i) follows from the corollary to Proposition~\ref{prop:II.2.4.1}.

\runhead{Exercise}

Show that the characters of $S_m$ correspond in a one-one way to the
homomorphisms $\chi\colon I \to \bar{\Q}^*$ having the following two
properties:
\begin{alist}
\item $\chi(x) = 1$ if $x \in U_m$
\item For each embedding $\sigma$ of $K$ into $\bar{\Q}$, there exists
an integral number $n(\sigma)$ such that
\[
\chi(x) = \prod_{\sigma \in \Gamma} \sigma(x)^{n(\sigma)}
\]
for all $x \in K^*$.
\end{alist}

\subsechead{2.5}{\textit{l}-adic representations associated to a linear representation of $S_m$}

\runhead{1)  The \textit{l}-adic case}

Let $V_\ell$ be a finite-dimensional $\Q_\ell$-vector space and
\[
\phi\colon S_{m/\Q_\ell} \longrightarrow \GL_{V_\ell}
\]
a linear representation of $S_{m/\Q_\ell}$ in $V_\ell$.  This defines a
homomorphism
\[
\phi\colon S_m(\Q_\ell) \longrightarrow \GL_{V_\ell}(\Q_\ell)
= \Aut(V_\ell)
\]
which is continuous for the $\ell$-adic topologies of those groups.

By composition with the map
$\epsilon_\ell\colon G^{\mathrm{ab}} \to S_m(\Q_\ell)$ defined in
\S\ref{sec:II.2.3}, we get a map
\[
\phi_\ell = \phi \circ \epsilon_\ell\colon G^{\mathrm{ab}}
\longrightarrow \Aut(V_\ell),
\]
i.e.  an \emph{abelian} \textit{l}-adic representation of $K$ in
$V_\ell$.

%% NOTE (II-18): "$\Sigma_k$" below is AS PRINTED (lowercase $k$, unlike
%% $\Sigma_K$ everywhere else); kept as printed.
\begin{proposition}
\begin{alist}[label=\alph*)]
\item The representation $\phi_\ell$ is semi-simple.
\item Let $v \in \Sigma_k$, with $v \notin \Supp(m)$ and
$p_v \neq \ell$.  Then $\phi_\ell$ is unramified at $v$; the
corresponding Frobenius element $F_{v,\phi_\ell} \in \Aut(V_\ell)$ is
equal to $\phi(F_v)$, where $F_v$ denotes the element of $S_m(\Q)$
defined in \S\ref{sec:II.2.3}.
\item The representation $\phi_\ell$ is rational
(chap.~\ref{chap:I}, \S\ref{sec:I.2.3}) if and only if $\phi$ can be
defined over $\Q$ (cf.~\S\ref{sec:II.2.4}).
\end{alist}
\end{proposition}

Since $S_m$ is a group of multiplicative type, all its representations
can be brought to diagonal form on a suitable extension of the ground
field; hence \emph{a}).  Assertion \emph{b}) follows from
\S\ref{sec:II.2.3}, and assertion \emph{c}) follows from
Proposition~\ref{prop:II.2.4.2} of \S\ref{sec:II.2.4}.

\runhead{Remark}

Let us identify $\phi_\ell$ with the corresponding homomorphism of the
idèle group $I$ into $\Aut(V_\ell)$.  Then

d) $\Ker(\phi_\ell)$ \emph{contains} $U_{v,m}$ \emph{if}
$v \notin \Supp(m)$, $p_v \neq \ell$.

e) Let $\phi_T\colon T_{/\Q_\ell} \to \GL_{V_\ell}$ be defined by
composing $T_{/\Q_\ell} \to S_{m/\Q_\ell}$ with $\phi$.  If $x$ belongs
to the open subgroup $U_{\ell,m} = \prod\limits_{v|\ell} U_{v,m}$ of
$T(\Q_\ell)$, one has
\[
\phi_\ell(x) = \phi_T(x^{-1})\,.
\]

These properties follow readily from those of $\epsilon_\ell$.

\runhead{2)  The rational case}

Let now $V_o$ be a finite dimensional vector space over $\Q$ and
$\phi_o\colon S_m \to \GL_{V_o}$ a linear representation of $S_m$.  For
each prime number $\ell$ we may apply the preceding construction to the
representation $\phi_{o/\ell}\colon S_{m/\Q_\ell} \to \GL_{V_\ell}$,
where $V_\ell = V_o \tens \Q_\ell$; we then get an \textit{l}-adic
representation $\phi_\ell\colon G^{\mathrm{ab}} \to \Aut(V_\ell)$.

\begin{theorem}
\begin{nlist}
\item The $\phi_\ell$ form a strictly compatible system of rational
abelian semi-simple representations.  Its exceptional set is contained
in $\Supp(m)$.
\item For each $v \notin \Supp(m)$ the Frobenius element of $v$ with
respect to the system $(\phi_\ell)$ is the element $\phi_o(F_v)$ of
$\Aut(V_o)$.
\item There exist infinitely many primes $\ell$ such that $\phi_\ell$
is diagonalizable over $\Q_\ell$.
\end{nlist}
\end{theorem}

The first two assertions follow directly from the proposition above.
To prove the third one, note first that there exists a finite extension
$E$ of $\Q$ over which $\phi_o$ becomes diagonalizable.  If $\ell$ is a
prime number which splits completely in $E$, one can embed $E$ into
$\Q_\ell$ and this shows that $\phi_\ell$ is diagonalizable.  Assertion
3) now follows from the well-known fact that there exist infinitely
many such $\ell$ (this is, for instance, a consequence of
\v{C}ebotarev's theorem, cf.~chap.~\ref{chap:I}, \S\ref{sec:I.2.2}).

\runhead{Remark}

The Frobenius elements $\phi_o(F_v) \in \Aut(V_o)$ can also be defined
using the homomorphism
\[
\phi_o \circ \epsilon\colon I \longrightarrow S_m(\Q)
\longrightarrow \Aut(V_o).
\]

Note that their eigenvalues generate \emph{a finite extension} of $\Q$;
indeed they are contained in any field over which $\phi_o$ can be
brought in diagonal form.

\begin{exercises}
\begin{nlist}
\item Let $\phi_o\colon S_m \to \GL_{V_o}$ be a linear representation
of $S_m$, and let $\ell$ be a prime number.
\begin{alist}[label=\alph*)]
\item Show that the Zariski closure of $\Ima(\phi_\ell)$ is the
algebraic group $\phi_o(S_m)$.  (Use the fact that
$\Ima(\epsilon_\ell)$ is Zariski dense in $S_m$,
cf.~\S\ref{sec:II.2.3}.)
\item Let $\mathfrak{s}_m$ be the Lie algebra of $S_m$ and
$\phi_o(\mathfrak{s}_m)$ be its image by $\phi_o$, i.e.  the Lie algebra
of $\phi_o(S_m)$.  Show that the Lie algebra of the $\ell$-adic Lie
group $\Ima(\phi_\ell)$ is $\phi_o(\mathfrak{s}_m) \tens \Q_\ell$.
(Use the fact that $\Ima(\epsilon_\ell)$ is open in $S_m(\Q_\ell)$,
cf.~\S\ref{sec:II.2.3}.)
\end{alist}

\item \begin{alist}[label=\alph*)]
\item Show that there exists a unique one-dimensional representation
\[
\mathrm{N}\colon S_m \longrightarrow \Gm
\]
such that $\mathrm{N}(F_v) = \mathrm{N}v \in \Q^*$ for all
$v \notin \Supp(m)$.
\item Show that the morphism
$T \to S_m \xrightarrow{\ \mathrm{N}\ } \Gm$ is the one induced by the
norm map from $K$ to $\Q$.
\item Show that the \textit{l}-adic representation defined by
$\mathrm{N}$ is isomorphic to the representation $V_\ell(\mu)$ defined
in chap.~\ref{chap:I}, \S\ref{sec:I.1.2}.
\end{alist}
\end{nlist}
\end{exercises}

\subsechead{2.6}{Alternative construction}

Let $\phi_o\colon S_m \to \GL_{V_o}$ be as in \S\ref{sec:II.2.5}.  If we
compose $\phi_o$ with the map $\epsilon\colon I \to S_m(\Q)$ defined in
\S\ref{sec:II.2.2}, we obtain a homomorphism
\[
\phi_o \circ \epsilon\colon I \longrightarrow \GL_{V_o}(\Q)
= \Aut(V_o).
\]

Conversely:

\begin{proposition}
Let $f\colon I \to \Aut(V_o)$ be a homomorphism.  There exists a
$\phi_o\colon S_m \to \GL_{V_o}$ such that $\phi_o \circ \epsilon = f$
if and only if the following conditions are satisfied:
\begin{alist}
\item The kernel of $f$ contains $U_m$
\item There exists an algebraic homomorphism
$\psi\colon T \to \GL_{V_o}$ such that $\psi(x) = f(x)$ for every
$x \in K^* = T(\Q)$.
\end{alist}
Moreover, such a $\phi_o$ is unique.
\end{proposition}

\begin{proof}
The necessity of the conditions (a) and (b) is trivial.  Conversely, if
$f$ has properties (a), (b), it defines a \emph{homomorphism}
$I/U_m \to \Aut(V_o)$.  On the other hand, since $f$ and $\psi$ agree on
$K^*$, the \emph{morphism} $\psi$ is equal to $1$ on
$E_m = K^* \cap U_m$, hence on its Zariski-closure $\bar{E}_m$.  This
means that $\psi$ factors through
\[
T \longrightarrow T_m \longrightarrow \GL_{V_o}\,.
\]

By the universal property of $S_m$ (cf.~\S\ref{sec:II.1.3} and
\S\ref{sec:II.2.2}), the maps $I/U_m \to \GL_{V_o}(\Q)$ and
$T_m \to \GL_{V_o}$ define an algebraic morphism
$\phi_o\colon S_m \to \GL_{V_o}$, and one checks easily that $\phi_o$
has the required properties, and is unique.
\end{proof}

\runhead{Remark}

Since $U_m$ is open, property (a) implies that $f$ is
\emph{continuous} with respect to the discrete topology of
$\Aut(V_o)$.  Conversely, any continuous homomorphism
$f\colon I \to \Aut(V_o)$ is trivial on some $U_m$; moreover, there is a
smallest such $m$; it is called the \emph{conductor} of $f$.

\runhead{Exercise}

Let $m$ be a modulus and let $V_o$ be a finite dimensional $\Q$-vector
space.  For each $v \notin \Supp(m)$ let $F_v$ be an element of
$\Aut(V_o)$.  Assume
\begin{alist}
\item The $F_v$'s commute pairwise.
\item There exists an algebraic morphism $\psi\colon T \to \GL_{V_o}$
such that $\psi(\alpha) = \prod F_v^{v(\alpha)}$ for $\alpha \in K^*$,
$\alpha \equiv 1 \pmod{m}$, and $\alpha > 0$ at each real place.
\end{alist}
Show that there exists an algebraic morphism
$\phi_o\colon S_m \to \GL_{V_o}$ for which the Frobenius elements are
equal to the $F_v$'s.

\subsechead{2.7}{The real case}

The preceding constructions are relative to a given prime number
$\ell$.  However, they have an archimedean analogue, as follows:

Let $\pi\colon T \to S_m$ be the canonical map defined in
\S\ref{sec:II.2.3}, and let
\[
\pi_\infty\colon T(\R) \longrightarrow S_m(\R)
\]
be the corresponding homomorphism of real Lie groups.  Since
$T(\R) = (K \tens \R)^* = \prod\limits_{v \in \Sigma_K^\infty} K_v^*$,
we can identify $T(\R)$ with a direct factor of the idèle group $I$.
Let $\mathrm{pr}_\infty$ be the projection on this factor; the map
\[
\alpha_\infty = \pi_\infty \circ \mathrm{pr}_\infty\colon I
\longrightarrow T(\R) \longrightarrow S_m(\R)
\]
is continuous, and one checks as in \S\ref{sec:II.2.3} that
$\alpha_\infty$ coincides with $\epsilon$ on $K^*$.  One may then define
a map
\[
\epsilon_\infty\colon I \longrightarrow S_m(\R)
\]
by
\[
\epsilon_\infty(a) = \epsilon(a)\,\alpha_\infty(a^{-1}).
\]

One has $\epsilon_\infty(a) = 1$ if $a \in K^*$, hence $\epsilon_\infty$
may be viewed as \emph{a homomorphism of the idèle class group}
$C = I/K^*$ \emph{into the real Lie group} $S_m(\R)$.

The main difference with the ``finite'' case is that $\epsilon_\infty$
\emph{is not trivial on the connected component of} $C$, hence has no
Galois group interpretation.

When one composes $\epsilon_\infty\colon C \to S_m(\R)$ with a complex
character $S_{m/\C} \to \mathbf{G}_{m/\C}$, one gets a homomorphism
$C \to \C^*$, i.e.  a \emph{Grössencharakter} of $K$, in the sense of
Hecke.  It is easily seen that the characters obtained in this way
coincide with the ``\emph{Grössencharakter} of type $(A_o)$'' of Weil
(cf.~\cite{taniyama}, \cite{weil-characters}), whose conductor divide
$m$.

\runhead{Exercise}

Let
\[
e\colon I \longrightarrow S_m(\R) \times \prod_\ell S_m(\Q_\ell)
\]
be the map defined by $\epsilon_\infty$ and the $\epsilon_\ell$'s.

\begin{alist}[label=\alph*)]
\item Show that the image of $e$ is contained in the subgroup
$S_m(\A)$ of $S_m(\R) \times \prod\limits_\ell S_m(\Q_\ell)$, where
$\A$ denotes the ring of adèles of $\Q$, and that
$e\colon I \to S_m(\A)$ is continuous (for the natural topology of the
adelized group $S_m(\A)$).
\item Let $\pi_{\A}\colon T(\A) \to S_m(\A)$ be the map defined by
$\pi\colon T \to S_m$.  Show that, if one identifies $T(\A)$ with $I$
in the obvious way, one has
\[
e(x) = \epsilon(x)\,\pi_{\A}(x^{-1})
\]
where $\epsilon\colon I \to S_m(\Q) \subset S_m(\A)$ is the map defined
in \S\ref{sec:II.2.3}.  [Note that this gives an alternate definition
of the $\epsilon_\ell$'s.]
\item Show that $e(I)$ is not open in $S_m(\A)$ if $C_m \neq \{1\}$.
\end{alist}

\subsechead{2.8}{An example: complex multiplication of abelian varieties}

(We give here only a brief sketch of the theory, with a few indications
on the proofs.  For more details, see \cite{shimura-taniyama},
\cite{taniyama}, \cite{weil-characters}, \cite{weil-cm} and
\cite{serre-tate-good}.)

Let $A$ be an abelian variety of dimension $d$ defined over $K$.  Let
$\End_K(A)$ be its ring of endomorphisms and put
\[
\End_K(A)_o = \End_K(A) \tens \Q.
\]

Let $E$ be a number field of degree $2d$, and
\[
i\colon E \longrightarrow \End_K(A)_o
\]
be an injection of $E$ into $\End_K(A)_o$.  The variety $A$ is then said
to have ``complex multiplication'' by $E$; in the terminology of
Shimura-Taniyama, it is a variety of ``type (CM)''.

Let $\ell$ be a prime integer and define $T_\ell(A)$ and
$V_\ell = T_\ell(A) \tens \Q_\ell$ as in chap.~\ref{chap:I},
\S\ref{sec:I.1.2}.  These are free modules over $\Z_\ell$ and
$\Q_\ell$, of rank $2d$.  The $\Q$-algebra $\End_K(A)_o$ acts on
$V_\ell$; hence the same is true for $E$, and, by linearity, for
$E_\ell = E \tens_\Q \Q_\ell$.  One proves easily:

\begin{lemma}
$V_\ell$ is a free $E_\ell$-module of rank 1.
\end{lemma}

Let $\rho_\ell\colon \Gal(\bar{K}/K) \to \Aut(V_\ell)$ be the
\textit{l}-adic representation defined by $A$.  If
$s \in \Gal(\bar{K}/K)$, it is clear that $\rho_\ell(s)$ commutes with
$E$, hence with $E_\ell$.  But the lemma above implies that the
commuting algebra of $E_\ell$ in $\End(V_\ell)$ is $E_\ell$ itself.
Hence, $\rho_\ell$ may be identified with a homomorphism
\[
\rho_\ell\colon \Gal(\bar{K}/K) \longrightarrow E_\ell^*.
\]

Let now $T_E$ be the $2d$-dimensional torus attached to $E$ (as $T$ is
attached to $K$), so that $T_E(\Q_\ell) = E_\ell^*$, and $\rho_\ell$
takes values in $T_E(\Q_\ell)$.

\begin{theorem}[1]\label{thm:II.2.8.1}
\begin{alist}
\item The system $(\rho_\ell)$ is a strictly compatible system of
rational \textit{l}-adic representations of $K$ with values in $T_E$
(in the sense of chap.~\ref{chap:I}, \S\ref{sec:I.2.4}).
\item There is a modulus $m$ and a morphism
\[
\phi\colon S_m \longrightarrow T_E
\]
such that $\rho_\ell$ is the image by $\phi$ of the canonical system
$(\epsilon_\ell)$ attached to $S_m$, cf.~\S\ref{sec:II.2.3}.
\end{alist}
\end{theorem}

Moreover, the restriction of $\phi$ to $T_m$ can be given explicitly:

Let $t$ be the tangent space at the origin of $A$.  It is a $K$-vector
space on which $E$ acts, i.e.  an \emph{$(E,K)$-bimodule}.  If we view
it as an $E$-vector space, the action of $K$ is given by a homomorphism
$j\colon K \to \End_E(t)$.  In particular, if $x \in K^*$,
$\det\nolimits_E j(x)$ is an element of $E^*$; the map
$\det\nolimits_E j\colon K^* \to E^*$ is clearly the restriction of an
algebraic morphism $\delta\colon T \to T_E$.

\begin{theorem}[2]\label{thm:II.2.8.2}
The map $\delta\colon T \to T_E$ coincides with the composition map
$T \to T_m \to S_m \xrightarrow{\ \phi\ } T_E$.
\end{theorem}

\runhead{Example}

If $A$ is an \emph{elliptic curve}, $E$ is an imaginary quadratic field,
and the action of $E$ on the one-dimensional $K$-vector space $t$
defines an embedding $E \to K$.  The map
$\det\nolimits_E j\colon K^* \to E^*$ is just the \emph{norm} relative
to this embedding.

\runhead{Indications on the proofs of Theorems 1 and 2}

Part (a) of Theorem~\ref{thm:II.2.8.1} is proved as follows: Let $S$
denote the finite set of $v \in \Sigma_K$ where $A$ has ``bad
reduction''.  If $v \notin S$, and $\ell \neq p_v$, one shows easily
that $\rho_\ell$ is unramified at $v$ (the converse is also true, see
\cite{serre-tate-good}); moreover the corresponding Frobenius element
$F_{v,\rho_\ell}$ may be identified with the Frobenius endomorphism
$F_v$ of the reduced variety $\tilde{A}_v$.  But $F_v$ commutes with $E$
in $\End(\tilde{A}_v)_o$ and the commuting algebra of $E$ in
$\End(\tilde{A}_v)_o$ is $E$ itself (cf.~\cite{shimura-taniyama},
p.~39).  Hence $F_v$ belongs to $E^* = T_E(\Q)$ and this implies (a).

Theorem~\ref{thm:II.2.8.2} and part (b) of
Theorem~\ref{thm:II.2.8.1} are less easy; they are proved, in a somewhat
different form in \cite{shimura-taniyama} (see also
\cite{serre-tate-good}).  Note that one could express them (as in
\S\ref{sec:II.2.6}) by saying that \emph{there exists a homomorphism}
$f\colon I \to E^*$ (where $I$ denotes, as usual, the group of idèles
of $K$) \emph{having the following properties}:
\begin{alist}
\item $f$ \emph{is trivial on} $U_m$, \emph{for some modulus} $m$
\emph{with support} $S$.
\item \emph{If} $v \notin S$, \emph{the image by} $f$ \emph{of a
uniformizing parameter at} $v$ \emph{is the Frobenius element}
$F_v \in E^*$.
\item \emph{If} $x \in K^*$ \emph{is a principal idèle}, one has
$f(x) = \det\nolimits_E j(x)$.
\end{alist}

This is essentially what is proved in \cite{shimura-taniyama}, p.~148,
formula (3), except that the result is expressed in terms of ideals
instead of idèles, and $\det\nolimits_E j(x)$ is written in a different
form, namely
\[
\text{``}\ \prod_\alpha \mathrm{N}_{K/K^*}(x)^{\psi_\alpha}\ \text{''}.
\]

\runhead{Remark}

Another possible way of proving Theorems~\ref{thm:II.2.8.1}
and~\ref{thm:II.2.8.2} is the following:

Let $\ell$ be a prime integer distinct from any of the $p_v$,
$v \in S$.  One then sees that the Galois-module $V_\ell$ is of
Hodge-Tate type in the sense of chap.~\ref{chap:III},
\S\ref{sec:III.1.2} (indeed, the corresponding local modules are
associated with \emph{\textit{l}-divisible groups}, and one may apply
Tate's theorem \cite{tate-pdiv}).  Hence $\rho_\ell$ is ``locally
algebraic'' (chap.~\ref{chap:III}, \emph{loc. cit.}), and using the
theorem of chap.~\ref{chap:III}, \S\ref{sec:III.2.3} one sees it defines
a morphism $\phi\colon S_m \to T_E$.  One has
$\phi \circ \epsilon_\ell = \rho_\ell$ by construction; the same is true
for any prime number $\ell'$, since $\phi \circ \epsilon_{\ell'}$ and
$\rho_{\ell'}$ have the same Frobenius elements for almost all $v$.
This proves part (b) of Theorem~\ref{thm:II.2.8.1}.  As for
Theorem~\ref{thm:II.2.8.2}, one uses the explicit form of the
Hodge-Tate decomposition of $V_\ell$, as given by \cite{tate-pdiv},
combined with the results of the Appendix to chap.~\ref{chap:III}.
%% NOTE: the three lines above ("are associated with ... chap. III.")
%% are printed on II-29 (PDF 78), just before the Sec. 3 heading;
%% transcribed here to complete Sec. 2.8 (ch2_s3.tex starts at Sec. 3).
